\subsection{Representación del Cromosoma}

Cada individuo es un cromosoma binario de longitud $n$ (número de ítems):
$X = (x_1, x_2, \ldots, x_n)$, donde $x_i = 1$ indica que el ítem $i$ está
seleccionado. En memoria de GPU la población se almacena en un arreglo plano 1D
de tipo \texttt{uint8\_t[pop\_size $\times$ n\_items]} siguiendo un diseño \emph{Array of Structures} (AoS) lógico con la disposición \texttt{genes[ind $\times$ n\_items + gene]}.


\subsection{Flujo General del Algoritmo}

\begin{enumerate}
  \item \textbf{Inicialización}: generación aleatoria de la población con cuRAND.
  \item \textbf{Evaluación}: cálculo del fitness para cada individuo (kernel).
  \item \textbf{Reproducción}:
    \begin{enumerate}[label=\alph*.]
      \item Selección por torneo (tamaño 3, configurable).
      \item Cruzamiento de un punto (tasa 0.7).
      \item Mutación uniforme bit-flip (tasa 0.04 por gen).
      \item Reparación post-mutación (incompatibilidades, dependencias, capacidades).
    \end{enumerate}
  \item \textbf{Elitismo}: el 5\,\% superior de la población ($k \approx 0.05 \times \text{pop\_size}$) pasa directamente a la nueva generación.
  \item \textbf{Búsqueda del mejor}: recorrido lineal en CPU sobre el arreglo de
    fitness, priorizando individuos factibles (\texttt{is\_valid}). El cromosoma del
    mejor individuo se mantiene en GPU y se descarga \emph{una sola vez} al final
    mediante \texttt{get\_best()}, evitando transferencias masivas D$\to$H.
  \item \textbf{Criterio de término}: número fijo de generaciones (300).
\end{enumerate}

\subsection{Variante 1: Secuencial en CPU (\texttt{sequential})}

Implementación de referencia en C++17 con un único hilo. Sirve como línea base para
calcular el \emph{speed-up} de las variantes CUDA. Los operadores siguen el flujo
estándar usando \texttt{std::mt19937} como generador de números aleatorios.

\subsection{Variante 2: CUDA Básico (\texttt{cuda\_basic})}

Paralelización directa de los operadores principales en GPU:

\begin{itemize}
  \item \textbf{\texttt{fitness\_kernel}}: un hilo por individuo. Cada hilo recorre
    sus $n$ genes, acumula valor/peso/volumen y evalúa las cuatro clases de violaciones.
    Usa memoria global para todos los datos de la instancia.
  \item \textbf{\texttt{reproduce\_kernel}}: un hilo por hijo. Fusiona selección por
    torneo, cruzamiento, mutación y reparación en un único lanzamiento. El estado
    cuRAND de cada hilo se almacena en \texttt{curandState[pop\_size]}.
  \item \textbf{Transferencia eficiente}: solo los escalares de fitness, penalización
    y validez ($\sim$10 bytes/ind) se transfieren D$\to$H cada generación. El
    cromosoma del mejor individuo ($n$ bytes) se descarga de forma diferida
    (\emph{lazy}) en \texttt{get\_best()} al finalizar la evolución, eliminando
    la transferencia masiva de la población completa.
  \item \textbf{Medición de tiempos}: cada kernel y cada transferencia se mide con
    pares \texttt{cudaEvent\_t}, acumulando totales separados para fitness,
    reproducción, H$\to$D y D$\to$H.
  \item Los datos de la instancia (valores, pesos, volúmenes, restricciones) se
    transfieren \emph{una sola vez} antes del bucle evolutivo y permanecen en GPU
    durante toda la ejecución.
\end{itemize}

\subsection{Variante 3: CUDA Optimizado (\texttt{cuda\_optimized})}

Extiende \texttt{CUDABasic} con las siguientes optimizaciones medibles:

\begin{enumerate}
  \item \textbf{Memoria constante} (\texttt{\_\_constant\_\_}): cuando $n \leq 3000$,
    valores, pesos, volúmenes e IDs de categoría se suben a memoria constante
    (\texttt{c\_values}, \texttt{c\_weights}, etc.), cuyo \emph{cache} L1 dedicado
    reduce la latencia de acceso para lecturas uniforme-broadcast dentro de un warp.
  \item \textbf{Evaluación intra-bloque con \emph{shared memory}} 
    (\texttt{fitness\_kernel\_opt}): un bloque por individuo
    (\texttt{blockIdx.x = ind}). Los hilos del bloque procesan los genes en paralelo
    y utilizan memoria compartida para una reducción local (\emph{block reduction})
    acumulando valor, peso y volumen. Esto garantiza lecturas 100\% coalescentes y
    maximiza el ancho de banda. Solo se activa cuando $n \leq 3000$ y la instancia
    no tiene restricciones complejas (categorías, incompatibilidades, dependencias).
  \item \textbf{Streams CUDA}: la evaluación de fitness y la reproducción se solapan
    usando dos streams (\texttt{stream\_eval}, \texttt{stream\_repro}) con memoria
    \emph{pinned} (\texttt{cudaHostRegister}), permitiendo transferencias asíncronas
    y ocultando parte de la latencia de los kernels.
  \item \textbf{Tamaño de bloque óptimo}: se utiliza
    \texttt{cudaOccupancyMaxPotentialBlockSize} para determinar el tamaño de bloque
    que maximiza la ocupancia de los SM para el \texttt{reproduce\_kernel}.
  \item \textbf{Control de divergencia de warps}: los kernels de selección y mutación
    minimizan ramas dependientes del ID de hilo, manteniendo threads del mismo warp
    en la misma ruta de ejecución el mayor tiempo posible.
  \item \textbf{Parámetros escalares en memoria constante} (\texttt{c\_params}):
    los parámetros \texttt{FitnessParams} (capacidades, pesos de penalización, etc.)
    se almacenan en memoria constante y se transmiten a todos los hilos sin consumir
    registros adicionales.
  \item \textbf{Eventos de temporización}: cada kernel y cada transferencia se mide
    con \texttt{cudaEvent\_t} en el stream correspondiente, permitiendo desglosar
    con precisión los tiempos de fitness, reproducción y transferencias.
\end{enumerate}
